% Chapitre 3 : Combinaison de classifieurs et explicabilité
% ═══════════════════════════════════════════════════════════════════════════════
\section{Combinaison de Classifieurs et Explicabilité}
% ═══════════════════════════════════════════════════════════════════════════════

\begin{quote}
\itshape
« Un modèle qui prédit bien mais que personne ne comprend est aussi dangereux qu'un modèle
qui ne prédit pas du tout. »
\end{quote}

\vspace{0.5em}

Ce chapitre aborde une tension centrale de l'apprentissage automatique moderne~: les méthodes
les plus performantes — appelées \textit{ensembles de classifieurs} ou \textit{méthodes d'ensemble}
— sont précisément les plus difficiles à interpréter. Nous commencerons par comprendre pourquoi
le regroupement de plusieurs modèles crée un problème dit de \og boîte noire \fg{}, puis nous
analyserons les sources mathématiques de cette complexité, avant d'explorer en détail les
stratégies permettant de réconcilier performance et transparence.

% ═══════════════════════════════════════════════════════════════════════════════
\subsection{Le Problème de la Boîte Noire dans les Ensembles}
% ═══════════════════════════════════════════════════════════════════════════════

\subsubsection{Qu'est-ce qu'une boîte noire en apprentissage automatique ?}

Un modèle de machine learning est dit \textbf{boîte noire} (black box) lorsque sa structure
interne est trop complexe pour qu'un être humain puisse suivre mentalement le chemin de la
décision. On lui fournit des données en entrée, on obtient une prédiction en sortie, mais le
raisonnement intermédiaire reste opaque.

\paragraph{Définition (Boîte noire (Black Box)).}
Un modèle $f : \mathcal{X} \rightarrow \mathcal{Y}$ est qualifié de \textbf{boîte noire}
lorsqu'il est impossible, en pratique, pour un observateur humain de reconstituer le processus
de décision complet à partir de la représentation interne du modèle, même en ayant accès
à tous ses paramètres.

L'opposition classique se fait entre les modèles \textit{interprétables par conception}
(un arbre de décision de profondeur 3, une régression linéaire) et les modèles \textit{opaques
par nature} (un réseau de neurones profond, une forêt aléatoire de 500 arbres).

\subsubsection{Pourquoi les ensembles de classifieurs deviennent-ils des boîtes noires ?}

Imaginons un étudiant qui doit décider si une image médicale montre une tumeur bénigne
ou maligne. S'il consulte un seul médecin, il peut suivre son raisonnement pas à pas.
S'il consulte simultanément 1\,000 médecins et combine leurs avis selon une règle mathématique
complexe, il devient impossible de retracer la logique de la décision finale.

C'est exactement ce qui se passe avec les méthodes d'ensemble. Elles combinent plusieurs
modèles — appelés \textit{classifieurs de base} ou \textit{apprenants faibles} — pour
produire une prédiction plus robuste. Il existe trois grandes familles~:

\begin{enumerate}
  \item \textbf{Bagging (Bootstrap Aggregating)~:} les modèles sont entraînés en parallèle sur
        des sous-échantillons aléatoires du jeu de données d'entraînement. Exemple canonique~:
        la \textit{Forêt Aléatoire} (Random Forest).

  \item \textbf{Boosting~:} les modèles sont entraînés séquentiellement, chaque nouveau modèle
        cherchant à corriger les erreurs du précédent. Exemples~: AdaBoost, Gradient Boosting,
        XGBoost, LightGBM.

  \item \textbf{Stacking (Stacked Generalization)~:} les prédictions des modèles de base sont
        utilisées comme entrées d'un modèle de second niveau (\textit{méta-apprenant}) qui
        apprend à les combiner de façon optimale.
\end{enumerate}

\paragraph{Remarque (Pourquoi l'opacité est-elle dangereuse ?).}
Dans des domaines critiques comme la médecine (diagnostic de cancer), la finance (octroi
de crédit) ou la justice (aide à la décision judiciaire), un modèle opaque est problématique
car~: (1) il est impossible d'identifier et corriger des biais discriminatoires~; (2) les
réglementations (comme le RGPD en Europe) exigent un \og droit à l'explication \fg{}~;
(3) les experts métiers ne peuvent pas valider la cohérence du raisonnement.

\subsubsection{L'effet cascade de l'agrégation}

La décision finale d'un ensemble de $M$ classifieurs $\{h_1, h_2, \ldots, h_M\}$ pour
un vote majoritaire simple s'écrit~:

\begin{equation}
  \hat{y}(\mathbf{x}) = \underset{c \in \mathcal{C}}{\arg\max}
  \sum_{m=1}^{M} \mathbf{1}\bigl[h_m(\mathbf{x}) = c\bigr]
  \label{eq:vote_majoritaire}
\end{equation}

où $\mathbf{1}[\cdot]$ est la fonction indicatrice et $\mathcal{C}$ l'ensemble des classes.

Dans le cas continu (agrégation de probabilités), on obtient~:

\begin{equation}
  \hat{y}(\mathbf{x}) = \frac{1}{M} \sum_{m=1}^{M} h_m(\mathbf{x})
  \label{eq:moyenne_ensemble}
\end{equation}

Si l'on cherche à expliquer pourquoi $\hat{y}(\mathbf{x}_0) = 1$ pour une instance particulière
$\mathbf{x}_0$, il faut comprendre pourquoi \textit{chacun} des $M$ classifieurs a voté dans ce sens,
et comment leurs votes individuels ont été agrégés. Avec $M = 500$ arbres ayant chacun des
centaines de nœuds de décision, cette tâche devient humainement impossible sans outils
algorithmiques dédiés.

% ═══════════════════════════════════════════════════════════════════════════════
\subsection{Complexité des Ensembles de Modèles}
% ═══════════════════════════════════════════════════════════════════════════════

\subsubsection{La dimension de Vapnik-Chervonenkis et l'explosion de la capacité}

Pour mesurer formellement pourquoi les ensembles sont si difficiles à interpréter,
il faut comprendre la notion de \textbf{dimension de Vapnik-Chervonenkis} (dimension VC),
qui quantifie la richesse de l'espace des hypothèses d'un algorithme.

\paragraph{Définition (Dimension de Vapnik-Chervonenkis).}
Soit $\mathcal{H}$ un espace d'hypothèses (ensemble de fonctions de classification). La
\textbf{dimension VC} de $\mathcal{H}$, notée $\mathrm{VC}(\mathcal{H})$, est le plus grand
entier $d$ tel qu'il existe un ensemble de $d$ points $\{\mathbf{x}_1, \ldots, \mathbf{x}_d\} \subset \mathcal{X}$
que $\mathcal{H}$ peut \textit{pulvériser}, c'est-à-dire réaliser toutes les $2^d$
affectations binaires possibles sur ces points.

\textbf{Exemple concret~:} Un arbre de décision de profondeur maximale $k$ sur $p$ variables
binaires possède une dimension VC de l'ordre de $O(k \cdot p)$. C'est une valeur finie et
contrôlée. En revanche, une forêt aléatoire de $M$ arbres de profondeur $k$ génère une
frontière de décision dont la complexité combinatoire est de l'ordre de~:

\begin{equation}
  \mathrm{VC}(\mathcal{H}_{\text{forêt}}) = O\!\left(M \cdot k \cdot p \cdot \log(M)\right)
  \label{eq:vc_foret}
\end{equation}

Cette explosion de la dimension VC a deux conséquences directes sur l'interprétabilité~:

\begin{enumerate}
  \item \textbf{Perte de simulabilité~:} Un humain ne peut plus simuler mentalement
        le processus décisionnel. Il faut au minimum $O(\mathrm{VC}(\mathcal{H}))$ exemples pour
        comprendre le comportement du modèle.

  \item \textbf{Perte de décomposabilité~:} Il devient impossible d'attribuer
        un coefficient fixe et indépendant à chaque variable d'entrée. La contribution
        d'une variable $x_j$ dépend des valeurs de toutes les autres variables, de façon
        non linéaire et locale.
\end{enumerate}

\subsubsection{La décomposition biais-variance et le compromis d'opacité}

La théorie statistique de l'apprentissage décompose l'erreur attendue d'un modèle $h$ en
trois termes~:

\begin{equation}
  \mathbb{E}\!\left[(y - h(\mathbf{x}))^2\right]
  = \underbrace{\mathrm{Biais}^2[h(\mathbf{x})]}_{\text{erreur systématique}}
  + \underbrace{\mathrm{Variance}[h(\mathbf{x})]}_{\text{instabilité}}
  + \underbrace{\sigma^2_\epsilon}_{\text{bruit irréductible}}
  \label{eq:biais_variance}
\end{equation}

Les méthodes d'ensemble exploitent cette décomposition de manière différente selon
leur type~:

\begin{itemize}
  \item Le \textbf{Bagging} réduit principalement la \textit{variance}~: plusieurs modèles
        entraînés sur des sous-échantillons différents produisent des erreurs décorrélées,
        dont la moyenne est plus stable.

  \item Le \textbf{Boosting} réduit principalement le \textit{biais}~: les modèles successifs
        corrigent les erreurs résiduelles accumulées.
\end{itemize}

\textbf{Démonstration de la réduction de variance par le Bagging~:}

Soit $M$ classifieurs de base $h_1, \ldots, h_M$, entraînés sur des sous-ensembles
bootstrap indépendants, chacun ayant une variance $\sigma^2$ et une covariance pairwise
$\rho \sigma^2$ (avec $\rho$ le coefficient de corrélation entre deux prédicteurs).
La variance de leur moyenne est~:

\begin{equation}
  \mathrm{Var}\!\left(\frac{1}{M}\sum_{m=1}^{M} h_m(\mathbf{x})\right)
  = \rho \sigma^2 + \frac{1-\rho}{M}\,\sigma^2
  \label{eq:variance_bagging}
\end{equation}

Lorsque $M \to \infty$, le second terme tend vers zéro et la variance résiduelle
est $\rho \sigma^2$~: elle est entièrement déterminée par la \textbf{corrélation entre les
classifieurs}. C'est pourquoi les Random Forests introduisent la sélection aléatoire
de variables~: pour maximiser la diversité (réduire $\rho$).

\paragraph{Remarque (Paradoxe de l'opacité).}
La diversité entre classifieurs est \textit{à la fois} la cause de la performance supérieure
\textit{et} la cause principale de l'opacité. En forçant les modèles à raisonner sur des
sous-espaces de variables différents, on rend impossible toute explication unifiée et
cohérente de la décision finale.

\subsubsection{Le théorème de Condorcet et ses implications pour l'explicabilité}

Le \textbf{théorème du jury de Condorcet} (1785) fournit la justification théorique
de l'agrégation par vote majoritaire~:

\paragraph{Théorème (Jury de Condorcet).}
Supposons $M$ votants indépendants, chacun prenant la bonne décision avec probabilité
$p > 0.5$. La probabilité que la majorité prenne la bonne décision est~:
\begin{equation}
  P_M = \sum_{k=\lceil M/2 \rceil}^{M} \binom{M}{k} p^k (1-p)^{M-k}
\end{equation}
et $P_M \to 1$ lorsque $M \to \infty$.

\textbf{Ce que cela implique pour l'explicabilité~:} La force de l'ensemble réside
précisément dans la \textit{faiblesse et la spécialisation locale} de chaque apprenant
de base. Un apprenant performant sur une sous-région de l'espace n'a pas besoin d'être
correct ailleurs~: les autres compensent. Cela signifie qu'il est vain de chercher
à expliquer la décision finale par un raisonnement global cohérent~: la décision émerge
de la \textit{combinaison de raisonnements locaux et partiellement contradictoires}.

\subsubsection{L'identité de Krogh-Vedelsby~: diversité et erreur d'ensemble}

L'identité de Krogh et Vedelsby (1995) formalise mathématiquement le rôle de la diversité
dans la qualité d'un ensemble. Soit $\bar{h}(\mathbf{x}) = \frac{1}{M}\sum_{m=1}^M h_m(\mathbf{x})$
l'agrégat des prédictions, l'erreur quadratique moyenne de l'ensemble vérifie~:

\begin{equation}
  \mathbb{E}_{\mathbf{x}}\!\left[(y - \bar{h}(\mathbf{x}))^2\right]
  = \underbrace{\frac{1}{M}\sum_{m=1}^{M} \mathbb{E}\!\left[(y - h_m(\mathbf{x}))^2\right]}_{\text{erreur moyenne des classifieurs}}
  - \underbrace{\frac{1}{M}\sum_{m=1}^{M} \mathbb{E}\!\left[(\bar{h}(\mathbf{x}) - h_m(\mathbf{x}))^2\right]}_{\text{ambiguïté (diversité)}}
  \label{eq:krogh_vedelsby}
\end{equation}

\textbf{Interprétation~:} L'erreur de l'ensemble est \textit{toujours strictement inférieure}
à la moyenne des erreurs individuelles, l'écart étant exactement la diversité des
prédictions. Plus les classifieurs désaccordent sur leurs prédictions (grande diversité),
plus l'ensemble gagne en précision. Mais ce désaccord est précisément ce qui rend impossible
toute explication unifiée~: des classifieurs qui s'accordent sont interprétables~;
des classifieurs qui se contredisent délibérément ne le sont pas.

\subsubsection{L'accumulation des résidus dans le Boosting}

Le Gradient Boosting construit son modèle de façon additive et séquentielle~:

\begin{equation}
  F_T(\mathbf{x}) = \sum_{t=1}^{T} \eta \cdot h_t(\mathbf{x})
  \label{eq:gradient_boosting}
\end{equation}

où $\eta$ est le taux d'apprentissage et $h_t$ est le $t$-ième arbre entraîné sur le
gradient de la perte à l'itération $t$~:

\begin{equation}
  h_t = \underset{h \in \mathcal{H}}{\arg\min}
  \sum_{i=1}^{n} \left[-g_i^{(t)}\right] \cdot h(\mathbf{x}_i),
  \quad g_i^{(t)} = \frac{\partial \ell(y_i, F_{t-1}(\mathbf{x}_i))}{\partial F_{t-1}(\mathbf{x}_i)}
  \label{eq:gradient_residuel}
\end{equation}

\textbf{Conséquence pour l'explicabilité~:} Chaque arbre $h_t$ ne modélise pas la cible
$y$, mais le \textit{gradient résiduel}, c'est-à-dire la direction dans laquelle le modèle
précédent s'est trompé. Expliquer pourquoi $F_T(\mathbf{x}_0) = 0.87$ implique de justifier
non seulement le signal principal capturé par $h_1$, mais également l'ensemble
des micro-corrections $h_2, h_3, \ldots, h_T$, dont certaines compensent des erreurs
sur des régions très locales du jeu de données. Sans régularisation stricte, le modèle
mémorise du bruit~: toute explication extraite refléterait alors des \textit{artefacts
statistiques locaux} plutôt que des relations causales généralisables.

% ═══════════════════════════════════════════════════════════════════════════════
\subsection{Stratégies d'Explicabilité des Ensembles}
% ═══════════════════════════════════════════════════════════════════════════════

On distingue deux grandes catégories de stratégies~: les méthodes \textbf{post-hoc},
appliquées après l'entraînement d'un modèle opaque, et les architectures
\textbf{intrinsèquement interprétables} (\textit{glass-box}), conçues dès le départ pour
être transparentes. Nous verrons également les approches hybrides issues du méta-apprentissage.

\subsubsection{Méthodes post-hoc : limites de LIME sur les ensembles arborescents}

\paragraph{Principe de LIME.}
\textbf{LIME} (Local Interpretable Model-Agnostic Explanations) tente d'expliquer la
prédiction d'un modèle complexe $f$ en ajustant localement un modèle simple $g$
(régression Ridge ou Lasso) sur des instances perturbées autour de l'instance cible $\mathbf{x}_0$~:

\begin{equation}
  g^* = \underset{g \in \mathcal{G}}{\arg\min}
  \sum_{\mathbf{x}' \in Z} \pi_{\mathbf{x}_0}(\mathbf{x}') \cdot \ell\bigl(f(\mathbf{x}'), g(\mathbf{x}')\bigr)
  + \Omega(g)
  \label{eq:lime}
\end{equation}

où $Z$ est l'ensemble des échantillons perturbés, $\pi_{\mathbf{x}_0}(\mathbf{x}')$ est un noyau
de pondération gaussien qui donne plus de poids aux points proches de $\mathbf{x}_0$,
et $\Omega(g)$ est une pénalité de complexité.

\begin{algorithm}[H]
\caption{LIME appliqué à un classifieur d'ensemble $f$}
\begin{algorithmic}[1]
\Require Instance cible $\mathbf{x}_0$, modèle $f$, taille d'échantillon $N$, noyau $\pi$
\Ensure Explication locale $g^*$ (régression Ridge)
\State Initialiser $Z \leftarrow \emptyset$
\For{$i = 1$ \textbf{à} $N$}
    \State Générer $\mathbf{x}'_i$ par perturbation gaussienne de $\mathbf{x}_0$
    \State Calculer $\hat{y}_i = f(\mathbf{x}'_i)$ \Comment{Appel au modèle boîte noire}
    \State Calculer $w_i = \pi_{\mathbf{x}_0}(\mathbf{x}'_i) = \exp\!\left(-\frac{\|\mathbf{x}_0 - \mathbf{x}'_i\|^2}{\sigma^2}\right)$
    \State $Z \leftarrow Z \cup \{(\mathbf{x}'_i, \hat{y}_i, w_i)\}$
\EndFor
\State Résoudre : $g^* \leftarrow \arg\min_{g} \sum_{i} w_i \cdot (f(\mathbf{x}'_i) - g(\mathbf{x}'_i))^2 + \lambda\|g\|_1$
\State \Return $g^*$ (coefficients = importances des variables)
\end{algorithmic}
\end{algorithm}

\paragraph{Pourquoi LIME échoue sur les ensembles arborescents.}

Les ensembles arborescents partitionnent l'espace des caractéristiques en
\textit{hyper-rectangles orthogonaux}, créant une fonction de décision en escalier avec de
nombreuses discontinuités abruptes. Lorsque LIME génère des perturbations gaussiennes
autour de $\mathbf{x}_0$, certains points perturbés tombent dans le même hyper-rectangle que
$\mathbf{x}_0$ (même prédiction), d'autres traversent une frontière et obtiennent une prédiction
radicalement différente.

\textbf{Conséquence~:} l'ajustement linéaire devient extrêmement instable. Deux
exécutions successives de LIME sur la même instance peuvent produire des importances
de variables totalement différentes. Cette \textbf{instabilité est rédhibitoire} dans les
contextes réglementaires où la reproductibilité de l'explication est une exigence légale
(diagnostic clinique, octroi de crédit).

\paragraph{Exemple (Instabilité de LIME sur une forêt aléatoire).}
Supposons une forêt qui classifie les demandes de crédit. Pour le client A, une première
exécution de LIME indique que le \textit{revenu mensuel} est la variable la plus importante~;
une seconde exécution (avec un autre tirage aléatoire) indique que c'est \textit{l'historique
de remboursement}. Aucune des deux n'est fiable~: la banque ne peut pas s'appuyer sur
ces résultats pour justifier un refus de crédit devant un régulateur.

\subsubsection{TreeSHAP : la solution exacte fondée sur la théorie des jeux}

\paragraph{Valeurs de Shapley et garanties d'équité.}

Les \textbf{valeurs de Shapley} (SHAP), issues de la théorie des jeux coopératifs, définissent
la contribution équitable de chaque variable $j$ à la prédiction $f(\mathbf{x}_0)$~:

\begin{equation}
  \phi_j(f, \mathbf{x}) = \sum_{S \subseteq \mathcal{F} \setminus \{j\}}
  \frac{|S|!\,(|\mathcal{F}| - |S| - 1)!}{|\mathcal{F}|!}
  \left[f_{S \cup \{j\}}(\mathbf{x}_{S \cup \{j\}}) - f_S(\mathbf{x}_S)\right]
  \label{eq:shapley}
\end{equation}

où $\mathcal{F}$ est l'ensemble de toutes les variables, et $f_S$ est le modèle restreint
aux variables de l'ensemble $S$ (les variables absentes étant marginalisées).

Ces valeurs satisfont quatre axiomes fondamentaux~:
\begin{enumerate}
  \item \textbf{Efficacité~:} $\sum_{j=1}^{p} \phi_j = f(\mathbf{x}) - \mathbb{E}[f(\mathbf{x})]$ (les contributions somment à la déviation par rapport à la moyenne).
  \item \textbf{Symétrie~:} si deux variables $i$ et $j$ contribuent identiquement dans toutes les coalitions, $\phi_i = \phi_j$.
  \item \textbf{Nullité~:} si une variable n'affecte jamais la prédiction, $\phi_j = 0$.
  \item \textbf{Additivité~:} pour un modèle additif $f = f_1 + f_2$, $\phi_j(f) = \phi_j(f_1) + \phi_j(f_2)$.
\end{enumerate}

\textbf{Problème calculatoire~:} L'équation \eqref{eq:shapley} somme sur $2^{|\mathcal{F}|}$
sous-ensembles. Avec $p = 30$ variables, cela représente plus d'un milliard de calculs.
L'approche agnostique \textbf{KernelSHAP} approxime ce calcul par échantillonnage, avec
une complexité en $O(2^M)$, ce qui est rédhibitoire pour des ensembles massifs.

\paragraph{L'algorithme TreeSHAP : complexité polynomiale exacte.}

\textbf{TreeSHAP} exploite la structure arborescente pour calculer les valeurs de Shapley
\textit{exactes} avec une complexité réduite à~:

\begin{equation}
  O(T \cdot L \cdot D^2)
  \label{eq:complexite_treeshap}
\end{equation}

où $T$ est le nombre d'arbres, $L$ le nombre maximal de feuilles par arbre, et $D$
la profondeur maximale.

\textbf{Idée clé~:} Plutôt que de traiter le modèle comme une boîte noire et d'appeler
$f(\cdot)$ pour chaque coalition, TreeSHAP navigue dans la structure des nœuds de chaque
arbre et maintient, via de la \textit{programmation dynamique}, le compte du poids de
chaque feuille pour chaque coalition possible.

\begin{algorithm}[H]
\caption{TreeSHAP — calcul exact pour un arbre unique $t$}
\begin{algorithmic}[1]
\Require Arbre $t$, instance $\mathbf{x}_0$
\Ensure Valeurs de Shapley $\phi_j$ pour chaque variable $j$
\State Initialiser $\phi_j \leftarrow 0$ pour tout $j$
\State $\textsc{PathExplainer}(\text{nœud racine}, \mathbf{x}_0, \text{fraction}=1, \text{zero\_fraction}=1, \text{one\_fraction}=1)$
\Procedure{PathExplainer}{nœud $v$, $\mathbf{x}_0$, $\text{pf}$, $\text{zf}$, $\text{of}$}
  \If{$v$ est une feuille}
      \State Mettre à jour $\phi_j$ via les poids accumulés le long du chemin
  \Else
      \State Soit $j$ la variable de division en $v$ et $s$ le seuil
      \State $h \leftarrow$ fraction des données passant à gauche sous $v$
      \State Appel récursif pour le sous-arbre gauche avec fractions mises à jour
      \State Appel récursif pour le sous-arbre droit avec fractions mises à jour
  \EndIf
\EndProcedure
\State \Return $\phi = (\phi_1, \ldots, \phi_p)$
\end{algorithmic}
\end{algorithm}

\begin{table}[H]
\centering
\caption{Comparaison des méthodes SHAP pour les ensembles arborescents}
\renewcommand{\arraystretch}{1.4}
\begin{tabular}{>{\bfseries}p{3.8cm} p{3cm} p{3cm} p{3.8cm}}
\toprule
Critère & KernelSHAP & TreeSHAP & Fast TreeSHAP v2 \\
\midrule
Agnosticisme & Total (boîte noire) & Spécifique (arbres) & Spécifique (arbres) \\
Complexité & $O(2^M)$ approx. & $O(T \cdot L \cdot D^2)$ exact & Pré-calcul optimisé \\
Colinéarité & Restreinte & Paramétrable & Paramétrable \\
Vitesse & Très lente & Rapide & Très rapide \\
Usage temps réel & Non & Oui & Oui \\
\bottomrule
\end{tabular}
\label{tab:comparaison_shap}
\end{table}

\subsubsection{SIRUS : extraction de règles stables depuis les forêts aléatoires}

\paragraph{Le problème de l'instabilité des règles classiques.}

Une approche intuitive pour expliquer une forêt aléatoire consiste à en extraire des
règles logiques de la forme $\text{SI } (x_1 > 3.5) \text{ ET } (x_2 \leq 0.8) \text{ ALORS } y = 1$.
Cependant, les extracteurs de règles traditionnels souffrent d'une \textbf{instabilité
pathologique}~: un léger changement dans le jeu d'entraînement (un sous-échantillonnage
bootstrap différent) modifie la racine des arbres, déclenchant une cascade de règles
entièrement différentes.

\textbf{SIRUS} (Stable and Interpretable RUle Set) a été spécifiquement conçu pour
garantir la \textit{stabilité mathématique} des règles extraites d'une Random Forest.

\paragraph{Architecture de SIRUS en quatre phases.}

\begin{algorithm}[H]
\caption{Algorithme SIRUS}
\begin{algorithmic}[1]
\Require Forêt $\mathcal{F}$ (10\,000+ arbres, profondeur max 2), seuil $p_0$, quantiles $q$
\Ensure Ensemble de règles $\mathcal{R}^*$ stable et interprétable

\State \textbf{Phase 1 — Discrétisation~:}
\ForEach{variable continue $x_j$}
    \State Calculer les $q$ quantiles empiriques $\{c_{j,1}, \ldots, c_{j,q}\}$
    \State Remplacer $x_j$ par sa version discrétisée $\tilde{x}_j \in \{1, \ldots, q\}$
\EndFor
\Comment{Vocabulaire de coupures partagé entre tous les arbres}

\State \textbf{Phase 2 — Extraction des chemins~:}
\ForEach{arbre $t_k \in \mathcal{F}$}
    \ForEach{chemin racine-feuille $\pi$ dans $t_k$}
        \State Extraire la règle propositionnelle $r_\pi$ associée à $\pi$
        \State Ajouter $r_\pi$ au pool $\mathcal{R}$
    \EndFor
\EndFor

\State \textbf{Phase 3 — Filtrage par fréquence~:}
\ForEach{règle $r \in \mathcal{R}$}
    \State $\hat{p}(r) \leftarrow$ fréquence d'apparition de $r$ dans $\mathcal{F}$
    \If{$\hat{p}(r) \geq p_0$}
        \State Conserver $r$ dans $\mathcal{R}_{\text{freq}}$
    \EndIf
\EndFor

\State \textbf{Phase 4 — Élagage et pondération~:}
\State Supprimer les règles redondantes (combinaisons linéaires de règles plus courtes)
\State Ajuster les poids finaux par régression Ridge~: $\mathbf{w}^* = \arg\min_{\mathbf{w}} \|\mathbf{y} - R\mathbf{w}\|^2 + \lambda\|\mathbf{w}\|^2$
\State $\mathcal{R}^* \leftarrow$ liste finale de $\sim$10 règles pondérées

\State \Return $\mathcal{R}^*$
\end{algorithmic}
\end{algorithm}

\paragraph{Garantie de stabilité via l'indice de Dice-Sørensen.}

La stabilité de SIRUS est mesurée par l'indice de Dice-Sørensen entre deux ensembles
de règles $\mathcal{R}^{(1)}$ et $\mathcal{R}^{(2)}$ générés à partir de deux tirages
indépendants du même jeu de données~:

\begin{equation}
  \mathrm{DSC}(\mathcal{R}^{(1)}, \mathcal{R}^{(2)}) = \frac{2\,|\mathcal{R}^{(1)} \cap \mathcal{R}^{(2)}|}{|\mathcal{R}^{(1)}| + |\mathcal{R}^{(2)}|}
  \label{eq:dice_sorensen}
\end{equation}

Les études empiriques montrent que SIRUS atteint un DSC supérieur à 0.95 (contre moins
de 0.5 pour les extracteurs classiques), tout en préservant une précision prédictive
asymptotiquement proche de la forêt aléatoire d'origine.

\paragraph{Exemple (SIRUS pour le diagnostic médical).}
Appliqué à un jeu de données de diagnostic cardiaque (variables~: âge, pression artérielle,
taux de cholestérol, fréquence cardiaque maximale), SIRUS extrait typiquement un ensemble
de 8 à 12 règles du type~:
\begin{itemize}
  \item \textit{SI âge $>$ 55 ET cholestérol $>$ 240 ALORS risque élevé (poids~: 0.31)}
  \item \textit{SI fréquence\_max $\leq$ 120 ET pression\_art $>$ 140 ALORS risque élevé (poids~: 0.24)}
\end{itemize}
Un cardiologue peut valider chacune de ces règles médicalement, chose impossible avec les
500 arbres de la Random Forest sous-jacente.

\subsubsection{Les architectures Glass-Box : transparence par conception}

\paragraph{Les Explainable Boosting Machines (EBM).}

L'\textbf{Explainable Boosting Machine} (EBM) constitue l'avancée paradigmatique majeure~:
un modèle aussi performant qu'un Gradient Boosting, mais intrinsèquement interprétable.
L'EBM s'inscrit dans le cadre des Modèles Additifs Généralisés avec Interactions (GA2M)
dont la fonction d'hypothèse est~:

\begin{equation}
  g\!\left(\mathbb{E}[y]\right) = \beta_0 + \sum_{j=1}^{d} f_j(x_j) + \sum_{i \neq j} f_{ij}(x_i, x_j)
  \label{eq:ebm}
\end{equation}

où $g$ est la fonction de lien (logit pour la classification), $f_j(x_j)$ capture l'effet
univarié de la variable $j$, et $f_{ij}(x_i, x_j)$ modélise les interactions paires
les plus significatives.

\textbf{La clé de l'interprétabilité~: le Boosting cyclique (round-robin).}
Contrairement au Gradient Boosting classique qui entraîne des arbres multivariés profonds,
l'EBM entraîne séquentiellement un micro-arbre de profondeur 1 (stump) pour une seule
variable à la fois~:

\begin{algorithm}[H]
\caption{Entraînement d'une Explainable Boosting Machine}
\begin{algorithmic}[1]
\Require Données $(X, \mathbf{y})$, taux $\eta$ très faible, nombre d'itérations $T$
\Ensure Fonctions de forme $\{f_j\}_{j=1}^{d}$
\State Initialiser $F(\mathbf{x}) \leftarrow \bar{y}$ (prédiction constante)
\State Initialiser $f_j \leftarrow 0$ pour tout $j$
\For{$t = 1$ \textbf{à} $T$}
    \State Calculer les résidus~: $r_i \leftarrow y_i - F(\mathbf{x}_i)$ pour tout $i$
    \ForEach{variable $j \in \{1, \ldots, d\}$}
        \State Entraîner un stump $h_j^{(t)}$ sur $(x_j, r)$ uniquement
        \Comment{Un seul axe à la fois}
        \State Mettre à jour~: $f_j \leftarrow f_j + \eta \cdot h_j^{(t)}$
        \State Mettre à jour les résidus~: $r_i \leftarrow r_i - \eta \cdot h_j^{(t)}(x_{ij})$
    \EndFor
\EndFor
\State \Return $\{f_j\}_{j=1}^{d}$
\end{algorithmic}
\end{algorithm}

\textbf{Pourquoi cela garantit l'interprétabilité~:} Après convergence, chaque fonction
$f_j(x_j)$ est une fonction univariée exacte, visualisable sous forme d'un graphique 2D.
La contribution de la variable $j$ à la prédiction d'une instance $\mathbf{x}_0$ est exactement
$f_j(x_{0j})$, sans aucune approximation.

\begin{table}[H]
\centering
\caption{Comparaison entre Gradient Boosting (XGBoost) et EBM}
\renewcommand{\arraystretch}{1.4}
\begin{tabular}{>{\bfseries}p{3.8cm} p{4.5cm} p{4.5cm}}
\toprule
Métrique & Gradient Boosting (XGBoost) & Explainable Boosting Machine \\
\midrule
Topologie & Arbres profonds multivariés & Stumps univariés (round-robin) \\
Interactions & Illimitées et enchevêtrées & Limitées aux paires significatives \\
Explicabilité & Post-hoc (TreeSHAP requis) & Intrinsèque (fonctions 2D/3D) \\
Fidélité de l'explication & Approximative & Exacte (100\% fidélité) \\
Performance (AUC) & Référence & Statistiquement indiscernable \\
\bottomrule
\end{tabular}
\label{tab:xgboost_vs_ebm}
\end{table}

\paragraph{Distillation de connaissances (Knowledge Distillation).}

Lorsqu'un modèle opaque préexistant ne peut être remplacé par un EBM, la
\textbf{distillation de connaissances} (Hinton et al., 2015) permet de transférer
le savoir du modèle opaque vers un modèle transparent.

L'idée clé consiste à substituer les étiquettes binaires $y \in \{0, 1\}$ par les
\textbf{probabilités molles} (soft labels) produites par le modèle enseignant (teacher)~:
au lieu d'apprendre $y = 1$, le modèle étudiant (student) apprend
$[p_0 = 0.15, p_1 = 0.85]$. Ces probabilités molles encapsulent les
\textit{connaissances sombres} (dark knowledge) du modèle~: sa confiance, la similarité
qu'il perçoit entre les classes, la complexité locale de la frontière de décision.

La fonction de perte de la distillation est~:

\begin{equation}
  \mathcal{L}_{\text{dist}} = (1-\alpha)\,\mathcal{L}_{\text{CE}}(y, \hat{y}_s)
  + \alpha\,T^2\,\mathcal{L}_{\text{KL}}\!\left(\sigma\!\left(\frac{z_t}{T}\right),\, \sigma\!\left(\frac{z_s}{T}\right)\right)
  \label{eq:distillation}
\end{equation}

où $z_t$ et $z_s$ sont les logits du teacher et du student, $T$ est la \textit{température}
(hyperparamètre contrôlant la douceur des distributions), $\sigma$ est la fonction softmax,
$\alpha$ est un coefficient de pondération, et $\mathcal{L}_{\text{KL}}$ est la divergence
de Kullback-Leibler.

\paragraph{Exemple (Distillation d'une forêt en un arbre décisionnel).}
Un hôpital dispose d'une Random Forest à 300 arbres pour prédire les réadmissions
hospitalières (AUC = 0.89). Elle ne peut être remplacée car elle est intégrée dans
un système critique. En appliquant la distillation~:
\begin{enumerate}
  \item La forêt génère des probabilités molles $[0.78, 0.22]$, $[0.34, 0.66]$, etc.
        pour chaque patient du jeu d'entraînement.
  \item Un arbre de décision de profondeur 5 est entraîné sur ces probabilités molles.
  \item L'arbre résultant atteint AUC = 0.86 (vs 0.89 de la forêt) mais est
        entièrement lisible par un médecin~: 5 niveaux de règles simples.
\end{enumerate}

\subsubsection{Sélection Dynamique de Classifieurs (DCS/DES) et explicabilité contextuelle}

\paragraph{Principe de la sélection dynamique.}

Les approches \textbf{DCS} (Dynamic Classifier Selection) et \textbf{DES} (Dynamic Ensemble
Selection) adoptent une philosophie radicalement différente~: plutôt que de combiner tous
les classifieurs pour chaque instance, elles identifient dynamiquement, pour chaque
nouvelle instance $\mathbf{x}_0$, les classifieurs localement les plus compétents.

L'algorithme définit une \textbf{Région de Compétence} (RoC) en projetant $\mathbf{x}_0$ dans
l'espace de validation DSEL et en identifiant ses $k$ plus proches voisins~:

\begin{equation}
  \mathrm{RoC}(\mathbf{x}_0) = \bigl\{\mathbf{x}_{(1)}, \ldots, \mathbf{x}_{(k)}\bigr\}
  \subset \mathcal{D}_{\text{val}}
  \label{eq:roc}
\end{equation}

La compétence d'un classifieur $h_m$ sur cette région est évaluée par son accuracy locale~:

\begin{equation}
  \mathrm{Comp}(h_m, \mathbf{x}_0) = \frac{1}{k} \sum_{i=1}^{k}
  \mathbf{1}\!\left[h_m(\mathbf{x}_{(i)}) = y_{(i)}\right]
  \label{eq:competence}
\end{equation}

\begin{algorithm}[H]
\caption{Overall Local Accuracy (OLA) — DCS avec explicabilité}
\begin{algorithmic}[1]
\Require Instance $\mathbf{x}_0$, pool de classifieurs $\{h_m\}_{m=1}^M$, DSEL, $k$
\Ensure Prédiction $\hat{y}$ et justification contextuelle
\State Trouver $\mathrm{RoC}(\mathbf{x}_0) = k\text{-NN}(\mathbf{x}_0, \mathcal{D}_{\text{val}})$
\ForEach{classifieur $h_m$}
    \State Calculer $\mathrm{Comp}(h_m, \mathbf{x}_0)$ via l'équation \eqref{eq:competence}
\EndFor
\State Sélectionner $h^* = \arg\max_m \mathrm{Comp}(h_m, \mathbf{x}_0)$
\State $\hat{y} \leftarrow h^*(\mathbf{x}_0)$
\State \Return $\hat{y}$, $h^*$ (expert sélectionné), $\mathrm{RoC}(\mathbf{x}_0)$ (justification)
\end{algorithmic}
\end{algorithm}

\paragraph{Avantages pour l'explicabilité.}

La sélection dynamique offre une propriété explicative unique~: si le pool de classifieurs
est composé de modèles interprétables (arbres peu profonds, régressions logistiques),
le système global hérite de cette interprétabilité \textit{par héritage}. Pour justifier
la décision sur $\mathbf{x}_0$, il suffit d'expliquer~:
\begin{enumerate}
  \item Quel expert $h^*$ a été sélectionné~;
  \item Pourquoi il a été sélectionné (ses $k$ voisins historiques dans RoC)~;
  \item Quel est le raisonnement de $h^*$ (si $h^*$ est interprétable).
\end{enumerate}
C'est ce qu'on appelle le \textbf{raisonnement par analogie} (Case-Based Reasoning)~:
la décision est justifiée par des exemples historiques similaires, ce qui est naturellement
compréhensible par un expert métier.

\subsubsection{XStacking : le méta-apprentissage guidé par l'explication}

\paragraph{L'opacité fondamentale du Stacking classique.}

Dans un Stacking conventionnel, les prédictions des modèles de base constituent les
seules entrées du méta-apprenant. Si l'on dispose de $K$ modèles de base, le méta-apprenant
opère dans un espace de dimension $K$~: $\mathbf{z} = (\hat{y}_1, \ldots, \hat{y}_K)$.
L'\textbf{espace des variables d'origine est totalement écrasé}~: une méthode XAI appliquée
au méta-apprenant peut seulement dire \textit{«~le modèle A a contribué à 40\%~»}, mais
pas \textit{pourquoi}, ni via quelle variable d'entrée.

\paragraph{Architecture XStacking.}

\textbf{XStacking} résout ce problème en enrichissant l'espace du méta-apprenant avec
les vecteurs d'importance des caractéristiques générés pour chaque modèle de base~:

\begin{equation}
  \mathbf{z}_{\text{XStacking}} = \bigl[\hat{y}_1, \ldots, \hat{y}_K,\;
  \phi(f_1(\mathbf{x})), \ldots, \phi(f_K(\mathbf{x}))\bigr]
  \label{eq:xstacking}
\end{equation}

où $\phi(f_k(\mathbf{x})) = (\phi_1^{(k)}, \ldots, \phi_p^{(k)})$ est le vecteur des valeurs
SHAP du modèle $k$ sur l'instance $\mathbf{x}$. L'espace d'apprentissage du méta-modèle passe
de dimension $K$ à dimension $K + K \times d = K(1+d)$ (ou $K \times d$ si les prédictions
brutes sont omises).

\begin{table}[H]
\centering
\caption{Comparaison Stacking conventionnel vs XStacking}
\renewcommand{\arraystretch}{1.4}
\begin{tabular}{>{\bfseries}p{4cm} p{4.5cm} p{4.5cm}}
\toprule
Propriété & Stacking Conventionnel & Framework XStacking \\
\midrule
Entrées Niveau 1 & Prédictions $(\hat{y}_1, \ldots, \hat{y}_K)$ & Vecteurs SHAP + prédictions \\
Dimension méta-espace & $K$ & $K(1 + d)$ \\
Lien variables d'origine & Coupé (out-of-fold) & Propagé nativement \\
Explicabilité résultante & Post-hoc, faible traçabilité & Intrinsèque, auditable \\
Bénéfice secondaire & Réduction de variance & Élimination biais par corrélations fallacieuses \\
\bottomrule
\end{tabular}
\label{tab:stacking_vs_xstacking}
\end{table}

% ═══════════════════════════════════════════════════════════════════════════════
\subsection{Compromis Performance–Compréhension}
% ═══════════════════════════════════════════════════════════════════════════════

\subsubsection{Le mythe du compromis strict et sa déconstruction empirique}

Pendant longtemps, l'apprentissage automatique a été dominé par l'axiome suivant~:
\textit{un modèle plus performant est nécessairement moins interprétable}. Cet axiome
est aujourd'hui empiriquement réfuté par plusieurs résultats convergents.

\paragraph{L'EBM comme preuve de concept.}

Des évaluations extensives sur des jeux de données critiques (prédiction de faillite
bancaire, diagnostic de l'infarctus du myocarde, prédiction de risques géologiques)
montrent que les EBM atteignent un AUC statistiquement indiscernable de celui de XGBoost~:

\begin{table}[H]
\centering
\caption{Comparaison AUC — EBM vs XGBoost sur des jeux de données réels}
\renewcommand{\arraystretch}{1.4}
\begin{tabular}{lccc}
\toprule
Jeu de données & XGBoost & EBM & Différence \\
\midrule
Défaut de crédit (Taiwan) & 0.779 & 0.778 & $-0.001$ \\
Diagnostic cardiaque (UCI) & 0.921 & 0.918 & $-0.003$ \\
Mortalité soins intensifs & 0.862 & 0.859 & $-0.003$ \\
Détection fraude bancaire & 0.974 & 0.971 & $-0.003$ \\
\bottomrule
\end{tabular}
\label{tab:auc_comparison}
\end{table}

Ces différences sont statistiquement non significatives et bien inférieures aux variations
dues aux hyperparamètres ou à la variabilité d'échantillonnage. Le compromis n'est donc
pas une loi absolue, mais un artefact des architectures d'apprentissage antérieures.

\subsubsection{L'effet Rashomon et la multiplicité des modèles équivalents}

\paragraph{Définition formelle de l'ensemble de Rashomon.}

Le phénomène de \textbf{multiplicité des modèles} (\textit{Model Multiplicity} ou
\textit{Rashomon Effect}), formalisé par Leo Breiman, constitue l'une des découvertes les
plus déstabilisantes de l'apprentissage moderne~: dans des espaces à haute dimension,
il existe souvent un très grand nombre de modèles statistiquement équivalents, mais
reposant sur des logiques décisionnelles totalement incompatibles.

\paragraph{Définition (Ensemble de Rashomon $\mathcal{R}(\epsilon)$).}
Soit $f^*$ le modèle optimal dans l'espace des hypothèses $\mathcal{F}$, de perte empirique
$\mathcal{L}(f^*)$. L'\textbf{ensemble de Rashomon} est défini par~:
\begin{equation}
  \mathcal{R}(\epsilon) = \bigl\{f \in \mathcal{F} \;\mid\; \mathcal{L}(f) \leq \mathcal{L}(f^*) + \epsilon\bigr\}
  \label{eq:rashomon}
\end{equation}
Tous les modèles de $\mathcal{R}(\epsilon)$ ont un pouvoir de généralisation pratiquement
identique, mais peuvent s'appuyer sur des ensembles de variables et des logiques de décision
radicalement différents.

\paragraph{Implication pour l'explicabilité~: l'incertitude interprétative.}

Appliquer SHAP ou un PDP (Partial Dependence Plot) sur le \textit{meilleur modèle unique}
génère un \textbf{biais cognitif dangereux}~: cela donne l'illusion que le profil causal
obtenu est \textit{la seule vérité} sous-jacente aux données. Or, des centaines d'autres
modèles de performances équivalentes pourraient produire des explications contradictoires.

\textbf{La solution~: le Rashomon PDP.} On agrège les Graphiques de Dépendance Partielle
(PDP) sur l'ensemble $\mathcal{R}(\epsilon)$. Cela produit~:

\begin{itemize}
  \item \textbf{Stabilité concordante~:} si l'effet d'une variable $x_j$ est positif
        sur la quasi-totalité des modèles de $\mathcal{R}(\epsilon)$, l'explication est
        certifiée robuste et indépendante du choix d'optimisation.

  \item \textbf{Désaccord interne~:} si les PDP sont très dispersés entre les modèles
        de $\mathcal{R}(\epsilon)$, le système signale explicitement qu'il n'existe
        \textit{pas de consensus} dans les données sur le rôle de cette variable. Toute
        conclusion causale serait une sur-interprétation.
\end{itemize}

\paragraph{Exemple (Rashomon PDP en oncologie).}
Un modèle de prédiction de survie au cancer du sein est entraîné sur des variables
biologiques (expression génique, âge, stade). Le Rashomon PDP révèle que pour
la variable \textit{âge}~:
\begin{itemize}
  \item 95\% des modèles de $\mathcal{R}(0.01)$ concordent~: l'âge augmente le risque
        au-delà de 60 ans. $\rightarrow$ Explication certifiée robuste, actionnable cliniquement.
  \item Pour la variable \textit{marqueur HER2}~: 50\% des modèles indiquent un effet
        protecteur, 50\% un effet délétère. $\rightarrow$ Aucune conclusion causale ne doit
        être tirée~; une étude clinique dédiée est nécessaire.
\end{itemize}

\subsubsection{Le Consensus XAI~: agréger les explications pour réduire leur variance}

\paragraph{La méthode NormEnsembleXAI.}

Tout comme l'apprentissage d'ensemble stabilise les prédictions, le
\textbf{Consensus XAI} stabilise les explications en agrégeant les vecteurs d'importance
produits par différentes méthodes ou différentes exécutions. Le framework
\textbf{NormEnsembleXAI} structure ce processus en trois étapes~:

\begin{algorithm}[H]
\caption{NormEnsembleXAI — agrégation d'explications plurielles}
\begin{algorithmic}[1]
\Require Matrices d'explications $\{\Phi^{(1)}, \ldots, \Phi^{(K)}\}$ de sources distinctes
\Ensure Explication consensuelle stable $\bar{\Phi}$

\State \textbf{Étape 1 — Normalisation par méthode~:}
\ForEach{matrice $\Phi^{(k)}$}
    \State Appliquer une standardisation robuste~:
    \State $\tilde{\phi}^{(k)}_j \leftarrow \dfrac{\phi^{(k)}_j - \mathrm{median}(\Phi^{(k)})}
    {\mathrm{P}_{99}(\Phi^{(k)}) - \mathrm{P}_{01}(\Phi^{(k)})}$
    \Comment{Insensible aux valeurs aberrantes}
\EndFor

\State \textbf{Étape 2 — Harmonisation ordinale~:}
\State Vérifier que la normalisation préserve les ordres de rang entre variables

\State \textbf{Étape 3 — Agrégation~:}
\State Calculer le score de fiabilité~: $s_k = \mathrm{corr}(\Phi^{(k)}, f(\mathbf{x}))$
\State $\bar{\Phi} \leftarrow \sum_{k=1}^{K} s_k \cdot \tilde{\Phi}^{(k)} \Big/ \sum_{k=1}^K s_k$

\State \Return $\bar{\Phi}$
\end{algorithmic}
\end{algorithm}

\paragraph{Explanation Consistency Training (ECT).}

L'\textbf{Explanation Consistency Training} (ECT) va encore plus loin~: au lieu d'agréger
les explications après l'entraînement, ECT intègre la \textit{cohérence explicative} directement
dans la fonction de perte lors de l'entraînement. La perte totale est~:

\begin{equation}
  \mathcal{L}_{\text{ECT}} = \mathcal{L}_{\text{pred}} + \lambda_{\text{cons}} \cdot
  \frac{1}{N}\sum_{i=1}^{N} \left\|\phi(f, \mathbf{x}_i) - \phi(f, \mathbf{x}_i + \epsilon_i)\right\|^2
  \label{eq:ect}
\end{equation}

Le second terme pénalise les modèles dont les explications changent trop fortement sous
de petites perturbations $\epsilon_i$ des données d'entrée. L'ECT force l'ensemble à
converger non seulement vers une \textit{exactitude unanime} (même prédiction), mais
surtout vers une \textit{cohérence causale} (même justification). Des études sur des
données médicales (MIMIC-IV) montrent des améliorations de cohérence supérieures à
120\% sans dégradation significative de la performance prédictive.

\subsubsection{Guide de décision~: quelle stratégie adopter ?}

\begin{table}[H]
\centering
\caption{Guide de sélection de la stratégie d'explicabilité selon le contexte}
\renewcommand{\arraystretch}{1.6}
\begin{tabular}{p{3.5cm} p{2.5cm} p{2.5cm} p{3.5cm}}
\toprule
\textbf{Contexte} & \textbf{Contrainte principale} & \textbf{Méthode recommandée} & \textbf{Justification} \\
\midrule
Nouveau projet, domaine critique & Performance + légalité & EBM (GA2M) & Interprétabilité native, AUC comparable \\
Modèle existant, irremplaçable & Compatibilité & TreeSHAP + Distillation & Analyse locale précise \\
Audit réglementaire & Stabilité & SIRUS & Règles stables (DSC $>$ 0.95) \\
Décision sur instances rares & Contexte local & DCS/DES & Expertise locale certifiée \\
Haute dimensionnalité, $p > 100$ & Cohérence globale & Consensus XAI + Rashomon & Certifie la robustesse causale \\
Stacking existant & Traçabilité & XStacking & Lien bout-en-bout restauré \\
\bottomrule
\end{tabular}
\label{tab:guide_decision}
\end{table}

\subsection{Conclusion}

Ce chapitre a mis en lumière que l'opacité des ensembles de classifieurs n'est pas
un simple défaut technique, mais la conséquence directe des mécanismes mathématiques
qui fondent leur supériorité~: explosion de la dimension VC, exploitation délibérée
de la diversité (identité de Krogh-Vedelsby), accumulation de corrections résiduelles
(Boosting).

Face à cette complexité, l'ingénierie de l'explicabilité a développé un spectre de
solutions complémentaires~:

\begin{enumerate}
  \item \textbf{Post-hoc analytique~:} TreeSHAP calcule des valeurs de Shapley exactes
        en temps polynomial, en exploitant la structure des arbres.

  \item \textbf{Distillation symbolique~:} SIRUS extrait des règles SI-ALORS stables
        depuis les forêts aléatoires, sans perte significative de précision.

  \item \textbf{Architectures glass-box~:} Les EBM montrent empiriquement qu'interprétabilité
        et performance peuvent coexister sans compromis significatif.

  \item \textbf{Explicabilité contextuelle~:} La sélection dynamique (DCS/DES) offre
        une justification par analogie, adaptée aux instances localement complexes.

  \item \textbf{Méta-apprentissage transparent~:} XStacking restaure la traçabilité
        bout-en-bout du Stacking en enrichissant le méta-espace avec les vecteurs SHAP.

  \item \textbf{Robustesse de l'explication~:} Le Consensus XAI et les ensembles de
        Rashomon quantifient l'incertitude interprétative et certifient la stabilité causale.
\end{enumerate}

La frontière entre performance et compréhension, longtemps perçue comme infranchissable,
se révèle être le résultat d'une ingénierie insuffisante plutôt qu'une limite théorique
absolue. L'avenir des systèmes critiques repose sur des architectures où explicabilité
et précision se renforcent mutuellement, comme l'illustrent XStacking et les EBM.
